\begin{table*}[t]
\centering
\footnotesize
\setlength{\tabcolsep}{4pt}
\renewcommand{\arraystretch}{1.12}
\caption{Same entity, different meaning in \textcolor{encolor}{English} and \textcolor{encolor}{Hindi} idioms. Entities are matched through a shared English anchor; $n$ is the number of idioms per language the summary rests on. Every meaning shown was re-checked against those idioms and unsupported claims removed (\S\ref{sec:analysis-entity}). Hindi and Arabic idioms are romanised with a literal gloss. }
\label{tab:entity-cases-enhi}
\begin{tabular}{@{}p{0.085\textwidth}p{0.335\textwidth}p{0.365\textwidth}p{0.15\textwidth}@{}}
\toprule
Entity ($n$ en/hi) & \textcolor{encolor}{English idioms} & \textcolor{encolor}{Hindi idioms} & Shared \\
\midrule
\textbf{medicine}\newline (10/38)
& \textcolor{encolor}{unpleasant treatment or consequences one deserves or inflicts on others; a necessary but unpleasant truth, duty, or experience; healing, relief, or emotional benefit}\newline \textit{a dose of one’s own medicine}; \textit{bitter medicine}; \textit{laughter is the best medicine}
& \textcolor{encolor}{a remedy or treatment for illness; a necessary but unpleasant cure or hardship; advice or practice that should be followed for benefit}\newline \textit{maut ko dava nahin} ``no medicine for death''; \textit{kadii aushadhi bin piye, mithe na tan ki taap} ``without bitter medicine, no relief''; \textit{dava se parhej bada} ``avoidance is greater than medicine''
& necessary but unpleasant remedy \\
\textbf{grain}\newline (8/152)
& \textcolor{encolor}{small amount or trace; something to be skeptical about; a natural standard or inclination}\newline \textit{a grain of truth}; \textit{grain of truth}; \textit{take something with a grain of salt}
& \textcolor{encolor}{food or sustenance; agricultural produce or crop yield; something valuable to store, use, or protect}\newline \textit{akal ya sukal, anaj nikal} ``whether famine or plenty, take out grain''; \textit{anaj khao par bij bachao} ``eat grain but save seed''; \textit{admi anaj ka kida hai} ``man is a grain-eater''
& food or sustenance \\
\textbf{kick}\newline (15/35)
& \textcolor{encolor}{a forceful push or blow; amusement or pleasure; a period of enthusiasm or preference}\newline \textit{a kick in the pants}; \textit{get a kick out of}; \textit{for kicks}
& \textcolor{encolor}{a blow or act of kicking; harm, insult, or rejection; forceful correction or punishment}\newline \textit{uthte laat baithte ghunsa} ``standing kick, sitting punch''; \textit{gire mein sab chaar laat maarte hain} ``the fallen get four kicks''; \textit{ghode ko laat, aadmi ko baat} ``horse with kick, man with words''
& blow or kick \\
\textbf{plate}\newline (11/27)
& \textcolor{encolor}{a load of tasks or responsibilities; a busy schedule or burden; something easily obtained or within reach}\newline \textit{a full plate — a full schedule; a lot to do.}; \textit{have too much on one's plate — to have too many tasks or responsibilities to deal with}; \textit{on a plate — Within easy reach; easily achieved}
& \textcolor{encolor}{a serving dish or meal portion; a person's share of food or sustenance; a symbol of fortune, status, or conduct in life}\newline \textit{apni bhari thali chhode, doosre ki jhooti pattal nihare} ``leave full plate; covet others' leftovers''; \textit{abhage ki thali mein chhed} ``unlucky person's plate has a hole''; \textit{thali phootne par thikara hath aata hai} ``when the plate breaks, shards remain''
& food or sustenance \\
\textbf{shadow}\newline (14/28)
& \textcolor{encolor}{weakness or diminished vitality; fearfulness or timidity; certainty or lack of doubt}\newline \textit{a shadow of itself}; \textit{afraid of one's own shadow}; \textit{beyond a shadow of a doubt}
& \textcolor{encolor}{shelter or shade; changing fortune or alternating pleasure and pain}\newline \textit{chhaya badi maya hai} ``shade or shelter is very important''; \textit{idhar ki chhaya udhar jati hai} ``shade changes from here to there''; \textit{kuen ki chhaya kuen mein rehti hai} ``the well's shade stays in the well''
& shelter or shade \\
\bottomrule
\end{tabular}
\end{table*}

\begin{table*}[t]
\centering
\footnotesize
\setlength{\tabcolsep}{4pt}
\renewcommand{\arraystretch}{1.12}
\caption{Same entity, different meaning in \textcolor{encolor}{English} and \textcolor{zhcolor}{Arabic} idioms. Entities are matched through a shared English anchor; $n$ is the number of idioms per language the summary rests on. Every meaning shown was re-checked against those idioms and unsupported claims removed (\S\ref{sec:analysis-entity}). Hindi and Arabic idioms are romanised with a literal gloss. }
\label{tab:entity-cases-enar}
\begin{tabular}{@{}p{0.085\textwidth}p{0.335\textwidth}p{0.365\textwidth}p{0.15\textwidth}@{}}
\toprule
Entity ($n$ en/ar) & \textcolor{encolor}{English idioms} & \textcolor{zhcolor}{Arabic idioms} & Shared \\
\midrule
\textbf{bear}\newline (13/9)
& \textcolor{encolor}{very strong or powerful; irritable or bad-tempered; difficult or troublesome}\newline \textit{gruff as a bear — gruff; curt and unsociable. (*Also: as \textasciitilde{}.)}; \textit{hungry as a bear — Cliché very hungry. (*Also: as \textasciitilde{}.)}; \textit{like a bear with a sore head — to be in a very bad mood or irritable}
& \textcolor{zhcolor}{dangerous and predatory; a hard-to-control bad influence; a feared enemy or threat}\newline \textit{ibn al-dhib ma yitrabbash} ``son of wolf not raised''; \textit{la yidrab al-dhib wa la yjuwwa' al-ghanam} ``neither beat wolf nor starve sheep''; \textit{al-naghjah al-ayyatah ma yakulsh ibnha al-dhib} ``the bleating ewe's son wolf''
& dangerous, threatening creature \\
\textbf{heat}\newline (16/9)
& \textcolor{encolor}{pressure or trouble; anger or conflict; intensity or heightened activity}\newline \textit{catch heat}; \textit{put the heat on}; \textit{take the heat}
& \textcolor{zhcolor}{a neighborhood or local quarter; people of the neighborhood; a rough local street setting}\newline \textit{ish inta fi al-hara ya mankhul bila tara} ``what are you in the neighborhood''; \textit{tikhanigni fi zaffah witishtilih maaya fi harah} ``you fight me openly, reconcile secretly''; \textit{zayy wlad al-harah zummarah tijmahum wi-asaayah tifarraqhum} ``like neighborhood boys, whistle gathers them''
& none \\
\textbf{palm}\newline (9/16)
& \textcolor{encolor}{hand or palm as a place for money; bribery or payment; deception in passing something off}\newline \textit{grease someone's palm}; \textit{cross someone’s palm with silver}; \textit{palm off}
& \textcolor{zhcolor}{hand or palm as a measure of possession or small amount; ability or power to hold or obtain; a part of the body contrasted with other body parts}\newline \textit{jarada fi al-kaf wa-la alf fi al-hawa} ``a grasshopper in the hand''; \textit{wish bashush wa-la jawhar bimlu al-kaf} ``a smiling face is better''; \textit{al-tul 'a al-nakhl w-al-tukhn 'a al-gimmez} ``height is for palms''
& hand as a small holding place \\
\textbf{wool}\newline (10/22)
& \textcolor{encolor}{being genuine, solid, or of good quality; being firmly fixed in beliefs or character; being deceptive or misleading}\newline \textit{all wool and a yard wide}; \textit{dyed in the wool}; \textit{pull the wool over someone's eyes}
& \textcolor{zhcolor}{something of value or usefulness; something that returns to its original state; something soft, light, or easily stretched; something associated with deception or mistaken expectations}\newline \textit{rigi'i al-ghazl suf} ``the spun thread returned wool''; \textit{il-kalam zayy habl il-suf, kul ma tishidduh yitmatt} ``speech like a wool rope stretches''; \textit{huwwa hilit illi yijizz il-kalb suf} ``what can shear a dog wool''
& deception or misleading \\
\textbf{milk}\newline (9/18)
& \textcolor{encolor}{something wasted or irrecoverable; kindness or generosity; something used up or exploited for gain}\newline \textit{cry over spilled milk}; \textit{spilt milk}; \textit{the milk of human kindness}
& \textcolor{zhcolor}{generosity and hospitality; something mixed or confused; a relationship or situation that is easily spoiled or fragile}\newline \textit{sharru al-labani al-waliju} ``best milk is what enters''; \textit{samak laban tamr hindi} ``fish milk tamarind indian''; \textit{al-shirk zayy al-laban aqallaha haja ti'akkuruh} ``company like milk, slightest spoils it''
& something easily spoiled \\
\bottomrule
\end{tabular}
\end{table*}

\begin{table*}[t]
\centering
\footnotesize
\setlength{\tabcolsep}{4pt}
\renewcommand{\arraystretch}{1.12}
\caption{Same entity, different meaning in \textcolor{zhcolor}{Chinese} and \textcolor{encolor}{Hindi} idioms. Entities are matched through a shared English anchor; $n$ is the number of idioms per language the summary rests on. Every meaning shown was re-checked against those idioms and unsupported claims removed (\S\ref{sec:analysis-entity}). Hindi and Arabic idioms are romanised with a literal gloss. Chinese idioms are given with a literal English gloss. }
\label{tab:entity-cases-zhhi}
\begin{tabular}{@{}p{0.085\textwidth}p{0.335\textwidth}p{0.365\textwidth}p{0.15\textwidth}@{}}
\toprule
Entity ($n$ zh/hi) & \textcolor{zhcolor}{Chinese idioms} & \textcolor{encolor}{Hindi idioms} & Shared \\
\midrule
\textbf{shadow}\newline (79/28)
& \textcolor{zhcolor}{concealment or hiding; loneliness or being alone; close attachment or being inseparable}\newline \zh{匿影藏形} ``hide shadow and form''; \zh{孤形吊影} ``lonely form and shadow''; \zh{如影随形} ``like shadow following form''
& \textcolor{encolor}{shade as shelter or comfort; changing fortune or alternating conditions}\newline \textit{chhaya badi maya hai} ``shade is great illusion''; \textit{idhar ki chhaya udhar jati hai} ``shade moves this way there''; \textit{kuen ki chhaya kuen mein rehti hai} ``well's shade stays in well''
& shade as shelter or comfort \\
\textbf{god}\newline (145/378)
& \textcolor{zhcolor}{supernatural power or divine inspiration; a source of admiration or longing; a force associated with mental focus, calm, or mysterious change}\newline \zh{下笔如神} ``write like with divine power''; \zh{令人神往} ``make one greatly yearn for''; \zh{全神贯注} ``entire spirit fully focused''
& \textcolor{encolor}{god as protector of the helpless; god as provider and sustainer; god as the ultimate power behind fate and events}\newline \textit{andhe ka ishvar sahayak} ``god helps the blind''; \textit{akele-dukele ka allah beli} ``allah is companion of lone''; \textit{ajgar ke data ram} ``ram gives to the python''
& divine power and protection \\
\textbf{ant}\newline (33/60)
& \textcolor{zhcolor}{crowded massing together; petty or insignificant struggle; people drawn together in large numbers}\newline \zh{云屯蚁聚} ``cloud-gathered ant-assembled''; \zh{蚁斗蜗争} ``ant fights snail争''; \zh{蚁聚蜂屯} ``ants gather, bees mass''
& \textcolor{encolor}{smallness or insignificance; something that can be carried or received in tiny amount; a creature used in weather signs and comparisons}\newline \textit{aata hai hathi ke munh, jata hai chinti ke munh} ``comes into elephant's mouth, goes into ant's mouth''; \textit{kito ko kan, hathi ko man} ``ant gets a grain, elephant a maund''; \textit{chinti upar pahad} ``ant on top of mountain''
& smallness or insignificance \\
\textbf{hand}\newline (169/289)
& \textcolor{zhcolor}{ability to act or do something; control or authority over affairs; many people acting together}\newline \zh{一举手之劳} ``one hand's effort''; \zh{一手包办} ``one hand handles all''; \zh{一手遮天} ``one hand covers sky''
& \textcolor{encolor}{ability to do work oneself; control over one's own affairs or status; help or support from another person}\newline \textit{apna kaam aap bhala} ``one's own work is good''; \textit{apna haath jagannath} ``one's own hand jagannath''; \textit{apni izzat apne haath} ``one's honor in own hand''
& ability to act or do \\
\textbf{wind}\newline (411/32)
& \textcolor{zhcolor}{smooth progress without obstacles; changeable force or influence that can be followed or used; something that stirs up trouble or spreads things around}\newline \zh{一帆顺风} ``sail with favorable wind''; \zh{乘风破浪} ``ride wind break waves''; \zh{乘风转舵} ``ride wind turn rudder''
& \textcolor{encolor}{wind as a natural force affecting weather and crops; wind as something that causes trouble, disturbance, or change; wind as an unavoidable hardship or pressure}\newline \textit{uttar upaj bahu dhan dhan, khet baat sukh kare kisan} ``north wind brings rich crops''; \textit{kaun sa darakht hai jise hawa nahin lagi} ``which tree has not wind''; \textit{tez hawa se bachkar chalna padta hai} ``must walk avoiding strong wind''
& trouble or difficulty \\
\bottomrule
\end{tabular}
\end{table*}

\begin{table*}[t]
\centering
\footnotesize
\setlength{\tabcolsep}{4pt}
\renewcommand{\arraystretch}{1.12}
\caption{Same entity, different meaning in \textcolor{zhcolor}{Chinese} and \textcolor{zhcolor}{Arabic} idioms. Entities are matched through a shared English anchor; $n$ is the number of idioms per language the summary rests on. Every meaning shown was re-checked against those idioms and unsupported claims removed (\S\ref{sec:analysis-entity}). Hindi and Arabic idioms are romanised with a literal gloss. Chinese idioms are given with a literal English gloss. }
\label{tab:entity-cases-zhar}
\begin{tabular}{@{}p{0.085\textwidth}p{0.335\textwidth}p{0.365\textwidth}p{0.15\textwidth}@{}}
\toprule
Entity ($n$ zh/ar) & \textcolor{zhcolor}{Chinese idioms} & \textcolor{zhcolor}{Arabic idioms} & Shared \\
\midrule
\textbf{smoke}\newline (27/12)
& \textcolor{zhcolor}{disappearance or vanishing; something fleeting or passing quickly; abundance or multitude}\newline \zh{云飞烟灭} ``clouds fly, smoke vanishes''; \zh{烟云过眼} ``smoke and clouds pass by''; \zh{烟消云散} ``smoke disperses, clouds scatter''
& \textcolor{zhcolor}{a cause of harm or trouble; something used to force or pressure; a sign used in comparisons or proverbs}\newline \textit{ayyu fatan qatalahu ad-dukhan} ``what youth killed by smoke''; \textit{ad-dukhan al-qurayyib ya'mi} ``near smoke makes blind''; \textit{zayyi an-nahl ma yitla'ush illa ad-dukhan} ``bees emerge only with smoke''
& cause of disappearance or harm \\
\textbf{pearl}\newline (74/10)
& \textcolor{zhcolor}{something precious or valuable; elegant, polished speech or writing; something recovered or restored}\newline \zh{买椟还珠} ``buy casket return pearl''; \zh{口吐珠玑} ``mouth speak pearls and jades''; \zh{合浦珠还} ``hepu return pearl''
& \textcolor{zhcolor}{milk or a milk flow; something promised or expected but not delivered; a burden or troublesome rival}\newline \textit{razamatan wa la dirrah} ``a bleat and no milk''; \textit{la khayra fi razamatin la dirra ma'aha} ``no good in bleat without milk''; \textit{yamna'u darrahu wa darra ghayrihi} ``he withholds his milk and others'''
& none \\
\textbf{soul}\newline (105/26)
& \textcolor{zhcolor}{fear, panic, and loss of composure; something that can be seized, driven out, or separated from the body; a force that can strongly attract or shock}\newline \zh{丢魂丢魄} ``lose soul and spirit''; \zh{丧魂失魄} ``lose soul and soul''; \zh{勾魂摄魄} ``hook soul and seize spirit''
& \textcolor{zhcolor}{the self or inner person; a person’s feelings, desires, and resolve; something one can address, restrain, or blame}\newline \textit{tahammadi ya nafsu la hamida laki} ``praise yourself, no praiser''; \textit{al-hakimu yaqda'u al-nafs bil-kafaf} ``the wise restrains the self''; \textit{al-nafsu mulawwa'atun bihubbi al-'ajil} ``the self loves the immediate''
& inner self or spirit \\
\textbf{god}\newline (145/126)
& \textcolor{zhcolor}{divine or supernatural power; intense mental focus or vivid expression; trickery or deceptive behavior}\newline \zh{下笔如神} ``write like by divine aid''; \zh{全神贯注} ``whole spirit concentrated''; \zh{做神做鬼} ``act like gods and ghosts''
& \textcolor{zhcolor}{god as the one who acts and decrees; divine power in blessing or punishment; invoking god in curses or prayers}\newline \textit{amru allahu balghun yasadu bihi al-suda'aa'u wa yashqa bihi al-ashqiya'u} ``god's command reaches the happy''; \textit{bihamdi allahi la bihamdak} ``by god's praise, not yours''; \textit{abada allahu khadra'ahum} ``may god destroy their prosperity''
& supernatural power \\
\textbf{blessing}\newline (47/10)
& \textcolor{zhcolor}{good fortune or happiness; favorable outcome after misfortune; blessing as something to be shared or sought}\newline \zh{有福同享} ``share good fortune''; \zh{因祸得福} ``gain blessing from disaster''; \zh{祸兮福所倚，福兮祸所伏} ``fortune hides in misfortune''
& \textcolor{zhcolor}{blessing or baraka as abundance; blessing through cooperation and gathering; blessing as increased benefit from many people}\newline \textit{al-barakah fi kutr al-ayadi} ``blessing in many hands''; \textit{al-barakah fi al-lammah} ``blessing in the gathering''; \textit{al-barakah fi al-kuthrah} ``blessing in the multitude''
& good fortune and benefit \\
\bottomrule
\end{tabular}
\end{table*}

\begin{table*}[t]
\centering
\footnotesize
\setlength{\tabcolsep}{4pt}
\renewcommand{\arraystretch}{1.12}
\caption{Same entity, different meaning in \textcolor{encolor}{Hindi} and \textcolor{zhcolor}{Arabic} idioms. Entities are matched through a shared English anchor; $n$ is the number of idioms per language the summary rests on. Every meaning shown was re-checked against those idioms and unsupported claims removed (\S\ref{sec:analysis-entity}). Hindi and Arabic idioms are romanised with a literal gloss. }
\label{tab:entity-cases-hiar}
\begin{tabular}{@{}p{0.085\textwidth}p{0.335\textwidth}p{0.365\textwidth}p{0.15\textwidth}@{}}
\toprule
Entity ($n$ hi/ar) & \textcolor{encolor}{Hindi idioms} & \textcolor{zhcolor}{Arabic idioms} & Shared \\
\midrule
\textbf{wheat}\newline (61/10)
& \textcolor{encolor}{a staple crop tied to farming success and failure; something valuable, desirable, or the standard of quality; a basis for idioms about effort, planning, and deception}\newline \textit{aage gehun pichhe dhan, vaako kahiye bada kisan} ``front wheat, behind rice, great farmer''; \textit{aata nahin to dalia ho hi jayega} ``if no flour, porridge will happen''; \textit{gandumnuma jaufarosh} ``wheat-like barley seller''
& \textcolor{zhcolor}{a preferred or valued grain; something one has versus what others have; the expected end or proper destination of a thing}\newline \textit{ziwan baladna wa-la al-qamh al-salibi} ``our darnel, not the foreign wheat''; \textit{sha'irna wa-la qamh ghirna} ``our barley, not others' wheat''; \textit{al-qamh yidur w yiji al-tahun} ``the wheat turns and comes to mill''
& valued grain and standard of preference \\
\textbf{fool}\newline (53/14)
& \textcolor{encolor}{a person who is hard to teach or understand; someone whose actions or speech are useless or worthless; a person lacking wisdom or judgment}\newline \textit{aklamand ko ishaara, moorkh ko tamaacha} ``wise person, fool a slap''; \textit{aadhe ghade ka paani, moorkh ki kahi kahaani} ``half-pot water, fool's story''; \textit{gadha khilaaye paap na punn} ``feeding a donkey brings no merit''
& \textcolor{zhcolor}{a foolish or laughable person; someone who is self-centered or unconcerned with others; a person whose behavior is awkward or absurd}\newline \textit{qaalu: sabaah al-khair ya juha. qaal: dana lissa sarih} ``they said good morning, juha said''; \textit{qaalu: ya juha, imta tqum al-qiyamah? qaal: lamma amut ana} ``when will judgment day, when i die''; \textit{qaalu: ya juha, eih ahsan ayyak? qaal: lamma kunt a'abbi al-turab fi al-taqiyah} ``best days, when i filled cap''
& foolishness and absurdity \\
\textbf{lord}\newline (25/24)
& \textcolor{encolor}{a powerful or high-status person; a master or owner; a local notable whose respect depends on place and power}\newline \textit{apne-apne ghar sabhi thakur  apne ghar sabhi bade aur shaktishali hote hain. aashay yah hai ki apne ghar koi nahin dabta.} ``everyone is lord in own house''; \textit{chakar ko thakur bahut  chakar (naukar) ke liye malik bahut mil jaate hain. jo vyakti parishram karne wala hota hai, use malikoon ki koi kami nahin rehti.} ``servants find many masters''; \textit{thikane thakur pooja jaay  apne ilaqe mein hi thakur ki pooja hoti hai. jab koi sabal ya sampann vyakti apne kshetra se baahar kahin jaaye aur uska vahaan samman na ho, tab yah kaha jaata} ``lord is worshipped in place''
& \textcolor{zhcolor}{god; a lord/master who controls and arranges events; a higher power invoked for protection, relief, or justice}\newline \textit{ibn adam fi at-tafkir w-r-rabb fi at-tadbir} ``man thinks, lord arranges matters''; \textit{tibat nar tisbah ramad, laha rabb yudabbirha} ``fire becomes ashes; lord manages''; \textit{ya harib min qadaya, ma lak rabb siwaya} ``you have no lord but me''
& power and control \\
\textbf{grain}\newline (152/29)
& \textcolor{encolor}{food or sustenance; livelihood and economic value; something to be stored, saved, or used properly}\newline \textit{akal ya sukal, anaj nikaal} ``whether famine or plenty, keep grain''; \textit{anaj khao par beej bachao} ``eat grain but save seed''; \textit{aadmi anaj ka keeda hai} ``man depends on grain alone''
& \textcolor{zhcolor}{grain or seed as a small physical thing; food or crop; something tiny that can still matter or grow}\newline \textit{min al-habbati tansha'u al-shajaratu} ``from the grain grows the tree''; \textit{habbah tittaqall al-mizan} ``a grain makes the scale heavier''; \textit{al-hubb yatla' 'ala bidhrih} ``the crop comes from its seed''
& food or crop \\
\textbf{ox}\newline (147/8)
& \textcolor{encolor}{stubbornness or obstinacy; a hardworking farm animal used for plowing; something one can exploit for selfish gain}\newline \textit{aa bail mujhe maar} ``come ox, strike me''; \textit{aalasi ka kutta aur mehnati ka bail} ``lazy man's dog and hardworking ox''; \textit{apna bail kulhaadi naathe} ``one's own ox may be harnessed''
& \textcolor{zhcolor}{stubbornness or insistence; a farm animal used in practical work; a target of blame, mockery, or being made to bear others' actions}\newline \textit{huwa yaqoolu thawr wa al-akharu yaqoolu nahlibuhu} ``he says ox, the other says milk it''; \textit{naqool lahu thawr, yaqool ihlibhu} ``we say ox, he says milk it''; \textit{ghabar ya thawr 'ala qarnik} ``dust, ox, on your horn''
& stubbornness or obstinacy \\
\bottomrule
\end{tabular}
\end{table*}
